\chapter{Project Context and Scope}
\label{ch:context}

\section{Independent Research Project Setting}
\label{sec:setting}

This Final Year Project is an \emph{independent} academic research initiative carried
out under the scientific direction of Prof.\ Anuradha Kar, within the aivancity
Grande \'Ecole Programme PGE5 -- AI \& Data Science. The project is not embedded
within a commercial host company or a corporate product team. Instead, it operates as
a self-directed applied-research mission, producing open, reproducible scientific
artifacts: a positioned literature synthesis, a curated multi-jurisdictional compliance
dataset, an interactive analytics dashboard, and this Final Project Report.

This independent structure has direct methodological consequences. All data sources are
public---no proprietary clinical or commercial datasets are accessed---and all code,
schema definitions, and compiled reports are version-controlled for open reproducibility.
The absence of a commercial host organization is therefore a deliberate feature rather
than a limitation: it places scientific integrity and reproducibility as the primary
acceptance criteria.

\section{Scientific Direction and Supervision}
\label{sec:supervision}

The project is carried out under the sole direction of \textbf{Prof.\ Anuradha Kar},
who provides scientific oversight, methodological feedback, and critical evaluation of
the research design. Direction responsibilities include:

\begin{itemize}[noitemsep]
  \item validating the problem statement and research question formulation;
  \item reviewing the thematic literature synthesis and scoring ontology design;
  \item providing feedback on the RAG architecture and dashboard implementation
        choices;
  \item evaluating the written report against the aivancity PFE writing guidelines;
  \item approving the final submission for the oral defence jury.
\end{itemize}

There is no separate industry or internship supervisor for this project. All
scientific responsibility and quality governance flow through Prof.\ Anuradha Kar
and the project author.

\section{Role and Responsibilities of the Author}
\label{sec:role}

The project author combines the responsibilities of a research analyst, an AI/NLP
engineer, and a full-stack data-science developer:

\begin{itemize}[noitemsep]
  \item design the systematic review protocol and the seven-theme coding frame;
  \item curate 20 jurisdictional compliance profiles and theme scores with
        documented primary-source justification;
  \item implement the RAG module (corpus construction, hybrid retrieval, grounded
        generation, and the React and Streamlit interfaces);
  \item generate reproducible report artifacts: \LaTeX/PDF compilation, Python figure
        scripts, literature exports, and dataset versioning;
  \item maintain the open repository and ensure reproducibility for an independent
        examiner;
  \item iterate with Prof.\ Anuradha Kar on scientific positioning and writing quality.
\end{itemize}

Expected deliverables: Phase~1 topic framing; literature and policy synthesis; curated
compliance dataset; interactive dashboard (open at \repoLink{}); and this Final
Project Report.

\section{Mission Scope and Timeline}
\label{sec:mission}

The project progressed through successive milestones aligned with the programme
schedule:

\begin{itemize}[noitemsep]
  \item \textbf{Phase~1:} comparative review framing and dashboard concept;
  \item \textbf{Phase~2:} state-of-the-art drafting and dataset schema design;
  \item \textbf{Phase~3:} methodology formalization, dataset curation, and
        preliminary dashboard release;
  \item \textbf{Phase~4:} full-draft submission for supervisor review and
        integration of RAG module;
  \item \textbf{Phase~5:} final report, figures, code package, and defence
        preparation.
\end{itemize}

Throughout all phases, priority was given to a correct, reproducible MVP: open JSON
as the system of record, interactive comparison views, and a written report auditable
against primary sources. Fully supervised NLP classification and live legislative
monitoring were intentionally deferred to the roadmap rather than claimed as finished
deliverables.

\section{Data Accessed and Information Assets}
\label{sec:data}

All project data is drawn from publicly accessible regulatory and scientific sources:

\begin{itemize}[noitemsep]
  \item \textbf{Primary legal instruments:} FDA guidance and 21st Century Cures Act,
        EU AI Act 2024, EU MDR/IVDR, GDPR, MHRA AI guidance, NMPA and PMDA
        regulations, WHO ethics guidance~\cite{fda2021,euaiact2024,who2021};
  \item \textbf{International frameworks:} IMDRF AI/ML SaMD guidance, OECD AI
        Principles, GMLP;
  \item \textbf{Peer-reviewed literature:} 80+ references sourced from PubMed,
        Scopus, IEEE Xplore, and SSRN, covering AI regulation, SaMD compliance,
        ethics, and clinical validation;
  \item \textbf{Curated structured records:} \texttt{compliance\_dataset.json}
        containing 20 country objects with theme scores, narrative fields, legislation
        lists, and trend cards.
\end{itemize}

No identifiable patient-level or commercially sensitive data are processed. Data volume
is modest by machine-learning standards (tens of jurisdictional records rather than
millions of observations) but high in interpretive density: each country record encodes
narrative fields that must remain traceable to citations.

\section{Legal, Ethical, and Confidentiality Constraints}
\label{sec:legal}

\begin{itemize}[noitemsep]
  \item \textbf{GDPR / privacy:} only public law texts and published literature are
        processed; no special-category patient data in the current release;
  \item \textbf{Not legal advice:} compliance scores are comparative research
        indicators, not legal opinions;
  \item \textbf{Generative AI:} any AI-assisted drafting is declared in
        Appendix~\ref{app:genai}; analyses and conclusions remain the author's
        responsibility;
  \item \textbf{Conflicts of interest:} none declared; scores are not commissioned
        by any national regulator or vendor;
  \item \textbf{Open publication:} the code, dataset, and report cover only
        public-source analytics and are published openly on GitHub.
\end{itemize}

\section{Technical and Resource Constraints}
\label{sec:constraints}

Constraints that shaped design choices include:

\begin{itemize}[noitemsep]
  \item reliance on publicly documented device counts, which are heterogeneous across
        agencies and not directly comparable;
  \item English-dominant source coverage, introducing a potential language bias for
        non-Anglophone jurisdictions;
  \item local execution without paid regulatory-intelligence APIs, requiring all
        retrieval to be offline-capable;
  \item project timeline bounds that prioritize a robust curated MVP over fully
        supervised NLP classifiers.
\end{itemize}

Hardware requirements are laptop-class for both the Streamlit and React applications;
GPU resources are optional and reserved for future encoder fine-tuning. These
constraints directly justify: expert-curated scores rather than fully automated
extraction; versioned JSON rather than a multi-tenant database; and dual Streamlit
(prototyping) + React (presentation) implementations sharing a single data contract.

\section{Risk Register for the Mission}
\label{sec:risks}

\begin{table}[htbp]
\centering
\caption{Project risk register}
\label{tab:risks}
\begin{tabular}{@{}p{3.0cm}p{2.2cm}p{5.8cm}@{}}
\toprule
\textbf{Risk} & \textbf{Likelihood} & \textbf{Mitigation} \\
\midrule
Source obsolescence & Medium & Dataset \texttt{version}/\texttt{last\_updated} fields;
  event-driven update cadence \\
Misinterpretation of law & Medium & Non-legal-advice caveats; HITL score governance \\
Scope creep into clinical model building & Low & Explicit non-goals in
  Chapter~\ref{ch:intro} \\
Citation drift between report and code & Low & Single BibTeX/biblatex configuration
  shared by all report artifacts \\
UI maintenance cost & Low & Shared JSON system of record; clean schema versioning \\
\bottomrule
\end{tabular}
\end{table}

\section{Budgetary and Organizational Non-Goals}
\label{sec:nongoals}

The mission does not include paid access to commercial regulatory databases, clinical
trial sponsorship, or deployment inside a hospital EHR. The operating budget is
effectively zero beyond personal compute and open-source tooling. The project does not
replace institutional compliance functions; it produces a research artifact and
open-source platform that compliance practitioners could subsequently adapt.

\section{Supervision and Acceptance Criteria}
\label{sec:collab}

Under the direction of Prof.\ Anuradha Kar, the acceptance criteria for the Final
Project Report are:

\begin{enumerate}[noitemsep]
  \item a clearly articulated research question with testable scope;
  \item a thematically organised state of the art (not paper-by-paper enumeration);
  \item a reproducible scoring protocol with a published rubric and dataset version;
  \item a functional interactive dashboard launchable from a clean environment;
  \item critical discussion of construct, internal, and external validity threats;
  \item alignment with the aivancity PFE mandatory structure and IEEE citation
        practice.
\end{enumerate}

\section{Stakeholder Analysis}
\label{sec:stakeholders-context}

Even in an independent project setting, several stakeholder groups shape requirements:

\begin{table}[htbp]
\centering
\caption{Stakeholders and their relationship to the project}
\label{tab:stakeholders-ctx}
\begin{tabular}{@{}p{3.0cm}p{3.2cm}p{4.8cm}@{}}
\toprule
\textbf{Stakeholder} & \textbf{Role} & \textbf{Implication for the project} \\
\midrule
Prof.\ Anuradha Kar & Scientific director & Primary acceptance authority \\
aivancity jury & Defence evaluators & Structured report + oral defence \\
Peer researchers & Indirect users & Open-source reproducibility \\
Policymakers & Intended audience & Dashboard gap-identification views \\
MedTech developers & Intended audience & Country comparison + use-case readiness \\
Hospital compliance officers & Intended audience & Narrative drill-down views \\
\bottomrule
\end{tabular}
\end{table}

\section{Work Organization and Tools}
\label{sec:workorg}

The authoring stack combines \LaTeX{} for the report, Python/Streamlit for analytics,
TypeScript/React for the presentation layer, JSON for the dataset, and Git for version
control. The working rhythm alternates between (i)~source reading and score curation and
(ii)~writing, figure regeneration, and dashboard iteration. This mirrors scientific
software practice: data and narrative stay tightly coupled so that chapter claims
cannot silently drift from the repository state.

\section{Ethics Review Boundary}
\label{sec:ethicsboundary}

Because the project analyses public law and published secondary literature rather than
human-subject data, no clinical ethics-board protocol is required. The remaining ethical
risk is reputational: comparative scores could be misread as official rankings. Mitigation
strategies include explicit non-legal-advice caveats throughout the dashboard and report,
theme-level transparency, and a capacity-building framing in the conclusion that avoids
``naming and shaming'' interpretations.

\section{Summary of Contextual Constraints}
\label{sec:contextsummary}

In summary, this independent project operates under public-data inputs, zero cloud
budget, and a reproducibility mandate under the direction of Prof.\ Anuradha Kar.
Those constraints justify a curated comparative analytics MVP with an honest NLP
roadmap---prioritizing correctness and auditability over premature automation.
Chapter~\ref{ch:method} details the end-to-end architecture, search protocol,
seven-theme scoring ontology, and evaluation metrics that operationalize these choices.
